\section{Proofs of the path-level feasibility certificates}
\label{app:convex-feasibility}

\subsection{Convex observables and weighted kinetic energy}
\begin{proof}
Along an exact trajectory, two derivatives of $\phi(X_t)$ give
$u_t(X_t)^\top D^2\phi(X_t)u_t(X_t)+
\nabla\phi(X_t)\cdot\acc_t(X_t)$.
The global growth bounds and the finite second moment dominate these terms
when $D^2\phi$ is bounded.  Smooth compact cutoffs of $\phi$ have gradients
converging to $\nabla\phi$ in $L^2(p_t)$: both the gradient and the cutoff
error are bounded by $C(1+\|x\|)$, with the error supported in the tails.
Thus the acceleration condition removes the final integral in
\eqref{eq:convex-chain-rule}.  A convex Hessian makes the first integral
nonnegative.  Without that condition, the first integral is at least
$K_\phi(t)$, and the definition of $\varepsilon_t^{\rm acc}$ bounds the
absolute value of the final integral by
$\varepsilon_t^{\rm acc}\|\nabla\phi\|_{L^2(p_t)}$.  This proves
\eqref{eq:convex-residual-floor}.  The additional hypotheses stated for
unbounded tests justify the identical argument.
\end{proof}

The theorem implies the convex-order condition
\begin{equation}
 p_{(1-\lambda)s+\lambda t}
 \preceq_{\rm cx}(1-\lambda)p_s+\lambda p_t,
 \qquad 0\le s<t\le T,\quad 0\le\lambda\le1.
 \label{eq:convex-order-path}
\end{equation}
Here $\preceq_{\rm cx}$ compares expectations of convex functions, and the
right side is a mixture of endpoint measures.  To extend the proved
inequality beyond smooth tests, first mollify convex Lipschitz functions
and then approximate an arbitrary convex function from below by maxima of
finitely many supporting affine functions.  Subtracting one supporting
affine function permits monotone convergence; the mean is affine by
Corollary~\ref{cor:mean-action}.  This is a necessary condition, not an
existence theorem.  In particular, every straight stochastic interpolation
$\Law((1-t)X_0+tX_1)$ already satisfies the convex-observable test, regardless
of whether it has an acceleration-preserving deterministic generator.


For a smooth positive density, the weighted quantity $K_\phi$ has a linear
variational characterization.  If $B=D^2\phi$ satisfies $b_-I\preceq B\preceq b_+I$ for positive finite constants, a weighted Poincare
inequality and one finite-energy generating field give a solution $\chi$,
unique modulo constants in the weighted Sobolev completion of smooth
functions, of
\[
 \int \nabla\eta^\top B^{-1}\nabla\chi\,\dd p_t
 =\int\eta\,\partial_t p_t\,\dd x.
\]
Its velocity $v_\phi=B^{-1}\nabla\chi$ attains $K_\phi$.  Indeed, every other
admissible velocity differs by a zero-weighted-divergence field $r$, and
$\int r\cdot\nabla\chi\,\dd p_t=0$.  Expansion of the quadratic energy
proves minimality.  This permits a linear elliptic feasibility diagnostic
without asserting that the nonlinear correction problem is solvable.


\subsection{Covariance constraints and their limits}
\begin{proof}
Applying \eqref{eq:convex-chain-rule} to quadratic coordinate products,
or polarizing its quadratic-form identity, gives
$\frac{\dd^2}{\dd t^2}\E[X_tX_t^\top]=2\E[u_tu_t^\top]$.
Subtracting $m(t)m(t)^\top$ yields
$C''=2\E[(u_t-\xi)(u_t-\xi)^\top]$.
The displayed block matrix is the covariance of $(X_t,u_t(X_t))$.
For $Z=e\cdot(X_t-m(t))$ and $V=e\cdot(u_t(X_t)-\xi)$,
$q=\sigma_e^2$ satisfies $q'=2\E[ZV]$ and $q''=2\E[V^2]$.
Hence
\[
 \sigma_e''=
 \frac{\E[Z^2]\E[V^2]-(\E[ZV])^2}
 {(\E[Z^2])^{3/2}}\ge0
\]
by Cauchy--Schwarz.
\end{proof}

These tests form a strict hierarchy.  For example, a diagonal Gaussian
path with standard deviations $s_1(t)=1+t-t^2/10$ and
$s_2(t)=1+t+t^2$ on $[0,1/4]$ has positive $C''$ and passes the scalar
minimum-energy test, because
$M''/2-K_{\min}=s_1s_1''+s_2s_2''>0$.
It is nevertheless infeasible because $s_1''=-1/5<0$.
The next example also passes the full covariance test.


\subsection{The quartic obstruction and its residual floor}
\begin{proof}
The bounds $|\kappa|\le1$ and
$r|\kappa'(r)|\le5/(2e)<1$ give $\|Db\|_{\rm op}<2$.
Thus $G_t$ is a global smooth bi-Lipschitz map, with lower Lipschitz
constant at least $1-t^2\ge3/4$.  Its nonlinear terms have bounded
derivatives, and its inverse has uniformly bounded first derivatives on
the interval.  The generator and its material acceleration
$b\circ G_t^{-1}$ are globally Lipschitz.  The source is Gaussian and the
maps have linear growth, so every required moment is finite.  The generator
is radial, hence a gradient, with a potential of at most quadratic growth;
the weighted-gradient orthogonality argument establishes its minimum
kinetic energy.

For $R=\|Z\|$, $R^2/2$ is exponential with mean one, and
\[
 \E[R^{2\ell}e^{-jR^2/2}]
 =\frac{2^\ell\ell!}{(1+j)^{\ell+1}}.
\]
It follows exactly that
\[
 M(t)=2+\frac{t^2}{8}+\frac{23t^4}{576},\qquad
 K_{\min}(t)=\frac{23t^2}{144},\qquad
 M''-2K_{\min}=\frac14+\frac{23t^2}{144}>0.
\]
The mean is zero and $C(t)=q(t)I_2$, where
$q(t)=1+t^2/16+23t^4/1152$.
For $q=1+\alpha t^2+\beta t^4$ with $\alpha,\beta>0$,
\[
 qq''-\tfrac12(q')^2
 =2\alpha+12\beta t^2+6\alpha\beta t^4+4\beta^2t^6>0.
\]
Thus $\sqrt q$ is convex and the choice $B_t=(q'/2)I_2$ satisfies
\eqref{eq:covariance-lmi}.

For the convex observable $\phi(x)=\|x\|^4$, direct expansion instead gives
\begin{align*}
 F_\phi(t)
 &=8-\frac32t^2+\frac{73}{144}t^4
       +\frac{29}{9216}t^6+\frac{179}{73728}t^8,\\
 F_\phi''(t)
 &=-3+\frac{73}{12}t^2+\frac{145}{1536}t^4
       +\frac{1253}{9216}t^6<0
       \qquad(0\le t\le1/2).
\end{align*}
For the last inequality, the expression is increasing in $t^2$ and its
value at $t=1/2$ is $-867715/589824<0$.
The polynomial test is justified by the Gaussian moments and global
linear growth.  Theorem~\ref{thm:convex-feasibility} therefore excludes
every acceleration-preserving generator of this path, not only the
displayed base field or a selected correction family.  Finally,
$F_\phi''(0)=-3$ and
$\|\nabla\phi\|_{L^2(p_0)}=4\sqrt{\E R^6}=16\sqrt3$.
Equation~\eqref{eq:convex-residual-floor} gives
\eqref{eq:quartic-floor}.  Corollary~\ref{cor:sharp-local-residual}
then gives the stated local Euler coefficient lower bound.  For the
displayed radial base field itself, $\acc_t(G_tz)=b(z)$ is a gradient
with linear growth, so its projection equals itself and
$\varepsilon_t^{\rm acc}=\|b\|_{L^2(\rho)}=\sqrt{23}/12$
throughout the interval; this is larger than the universal floor.
\end{proof}
